\chapter{نتیجه‌گیری و کارهای آتی}

\section{جمع‌بندی پژوهش}

تحلیل احساسات ضمنی با تحلیل متن‌هایی سروکار دارد که نگرش نویسنده را بدون اتکای مستقیم به واژه‌های احساسی آشکار منتقل می‌کنند. در چنین متن‌هایی، قطبیت ممکن است از پیامد یک رویداد، برآورده‌شدن یا نشدن انتظار، مقایسه، توصیه یا رابطۀ میان هدف و بافت استنباط شود. دشواری مسئله زمانی بیشتر می‌شود که یک جمله چند موجودیت داشته باشد یا گزارشی خنثی در کنار نشانه‌ای احساسی قرار گیرد. ازاین‌رو، سامانۀ تحلیل باید هم احساس بیان‌نشده را بازیابی کند و هم از نسبت‌دادن نگرش به هدفی نادرست یا تفسیر افراطی یک واقعیت خنثی بپرهیزد.

هدف این پژوهش طراحی و ارزیابی سامانه‌ای چندمنبعی برای تحلیل احساسات ضمنی هدف‌محور بود. سامانۀ پیشنهادی یک پیش‌بینی مستقیم را در کنار مسیر استدلال چندمرحله‌ای الهام‌گرفته از \lr{THOR} قرار داد. هنگام اختلاف این دو مسیر، بازبینی ساختاریافته نوع خطا، برچسب پیشنهادی و سطح اطمینان را تولید کرد و کنترل‌گر قاعده‌محور تصمیمی میانی ساخت. در مرحلۀ نهایی، سیاستی محافظت‌شده که آستانه‌های آن فقط با اعتبارسنجی‌های داخلی بخش آموزش تنظیم شده بود، از میان منابع مستقیم، استدلالی و تشخیصی منبع مناسب را انتخاب کرد. این تفکیک اهمیت دارد؛ زیرا خروجی کنترل‌گر پاسخ نهایی سامانه نیست و انتخاب نهایی نیز به برچسب‌های طلایی آزمون دسترسی ندارد.

ارزیابی اصلی با \mbox{\lr{Qwen3 8B}} و از طریق \lr{Ollama} روی آزمون کامل \lr{SCAPT} انجام شد. نتایج نشان دادند که افزودن مراحل استدلالی به‌تنهایی دقت را تضمین نمی‌کند. هر دو پیکربندی استدلالی بررسی‌شده از پیش‌بینی مستقیم ضعیف‌تر بودند و کنترل‌گر قاعده‌محور نیز به‌تنهایی از خط پایه عبور نکرد. بااین‌حال، منابع مختلف دقیقاً در نمونه‌های یکسان خطا نمی‌کردند. سیاست انتخاب منبع توانست از این مکمل‌بودن استفاده کند و عملکرد نهایی را در هر دو دامنه و هر سه کلاس بهبود دهد.

بنابراین، یافتۀ محوری پژوهش برتری مستقل یک زنجیرۀ استدلال یا یک پرامپت خاص نیست. نتیجۀ اصلی این است که در تحلیل احساسات ضمنی، تولید استدلال و انتخاب پاسخ نهایی باید دو مسئلۀ جدا در نظر گرفته شوند. استدلال چندمرحله‌ای می‌تواند شواهد و پاسخ‌های مکمل فراهم کند، اما استفاده از آن باید تحت کنترل سازوکاری باشد که از دادۀ آموزش می‌آموزد هر منبع در چه شرایطی قابل اعتمادتر است.

\section{پاسخ به پرسش‌های پژوهش}

بر پایۀ نتایج فصل پنجم، پاسخ پرسش‌های مطرح‌شده در فصل نخست به شرح زیر است:

\begin{enumerate}
    \item \textbf{تفاوت پیش‌بینی مستقیم و استدلال چندمرحله‌ای.}
    پیش‌بینی مستقیم مسیر کوتاه‌تر و کم‌هزینه‌تری فراهم کرد و در ارزیابی اصلی از هر دو نسخۀ استدلالی دقیق‌تر بود. در مقابل، مسیر \lr{THOR} تصمیم را به تشخیص جنبه، استنتاج نظر و تعیین قطبیت تقسیم کرد و ردپایی قابل بررسی ساخت. این ردپا برای تحلیل و تولید منبع مکمل مفید بود، اما افزایش تعداد مراحل باعث انتقال خطای جنبه یا نظر به برچسب نهایی نیز شد. پس پاسخ بخش دوم پرسش منفی است: استدلال چندمرحله‌ای به‌تنهایی همواره موجب بهبود نمی‌شود و در تنظیم آزمایشی این پژوهش از پیش‌بینی مستقیم ضعیف‌تر بود.

    \item \textbf{نقش بازبینی ساختاریافته و کنترل‌گر.}
    بازبینی ساختاریافته اختلاف میان دو مسیر را به اطلاعاتی قابل استفاده شامل نوع خطا، پیشنهاد اصلاح و اطمینان تبدیل کرد. این اطلاعات امکان داد تصمیم‌ها بر اساس رده‌های مشخصی مانند احساس ازدست‌رفته، بیش‌تفسیر خنثی و جابه‌جایی دامنۀ هدف بررسی شوند. بااین‌حال، خروجی کنترل‌گر قاعده‌محور هنوز از خط پایۀ مستقیم پایین‌تر بود و برچسب تشخیصی نیز در هیچ نمونۀ آزمون به‌عنوان منبع نهایی انتخاب نشد. بنابراین، بازبینی و کنترل‌گر اختلاف را ردیابی‌پذیرتر و قابل‌کنترل‌تر کردند، اما شواهد موجود برتری مستقل پاسخ تشخیصی یا قواعد ثابت کنترل‌گر را نشان نمی‌دهند. فایدۀ عملی اطلاعات تشخیصی در این پژوهش، کمک به ساخت پروفایل انتخاب میان منابع مستقیم و استدلالی بود.

    \item \textbf{اثر انتخاب منبع تنظیم‌شده با دادۀ آموزش.}
    پاسخ این پرسش در مجموعۀ اصلی مثبت است. سامانۀ نهایی بدون استفاده از برچسب طلایی آزمون، دقت را از \lr{\setpersianfont ۰/۶۷۸۷۳۳} به \lr{\setpersianfont ۰/۷۲۳۹۸۲} و امتیاز اف‌یک ماکرو را از \lr{\setpersianfont ۰/۶۷۴۰۷۵} به \lr{\setpersianfont ۰/۷۱۹۱۱۹} رساند. این سیاست ۲۶ خطای پیش‌بینی مستقیم را اصلاح کرد و در مقابل، ۶ پاسخ درست مستقیم را تضعیف کرد. حاصل خالص، ۲۰ پاسخ درست بیشتر بود. در نتیجه، انتخاب وابسته به الگوهای آموخته‌شده از بخش آموزش از اتکای ثابت به مسیر مستقیم یا استدلالی مناسب‌تر بود. این نتیجه به سیاست، داده و مدل بررسی‌شده مقید است و به معنای برتری هر انتخاب‌گر در هر مجموعۀ دیگری نیست.

    \item \textbf{الگوهای اصلی خطا.}
    تحلیل کیفی سه الگوی برجسته را نشان داد. در حالت نخست، مدل پیامد یا انتظار ضمنی را بازیابی نمی‌کند و احساس مثبت یا منفی را از دست می‌دهد. در حالت دوم، یک گزارۀ خبری یا خنثی بیش‌ازحد تفسیر می‌شود و قطبیتی بدون شاهد کافی شکل می‌گیرد. در حالت سوم، نگرش مربوط به یک موجودیت یا وضعیت کلی به هدف تعیین‌شده منتقل می‌شود. مسیر استدلالی می‌تواند الگوی نخست را اصلاح کند، اما در صورت تثبیت نادرست هدف یا جنبه، دو الگوی دیگر را تشدید می‌کند. همچنین، پیش‌بینی مستقیم و سامانۀ نهایی در ۱۱۶ نمونۀ آزمون هر دو نادرست بودند؛ ازاین‌رو، بخشی از خطاها با انتخاب بهتر میان منابع موجود رفع نمی‌شود.

    \item \textbf{اثر انتخاب مدل و پشتانۀ اجرا.}
    این مقایسه روی ۹۰ نمونۀ مشترک آزمون انجام شد. در آن، \lr{Gemini} در هر دو مسیر مستقیم و استدلالی امتیاز بالاتری از اجرای متناظر \lr{Qwen} داشت. بااین‌حال، مسیر \lr{THOR} برای هر دو مدل از پیش‌بینی مستقیم همان مدل ضعیف‌تر بود. انتخاب منبع برای \lr{Gemini} بهبودی محدود ایجاد کرد و برای \lr{Qwen} در این زیرمجموعه تغییری در اف‌یک ماکرو نداد. چون خانواده مدل، زیرساخت اجرا و سایر مشخصات نیز هم‌زمان متفاوت‌اند، این مشاهده فقط رابطۀ انتخاب مدل و الگوی خطای آن با میزان سود مراحل بعدی را در تنظیم حاضر نشان می‌دهد و اثر علّی ظرفیت مدل را اثبات نمی‌کند. با توجه به اندازۀ محدود آزمون مشترک، این نتیجه شاهدی تکمیلی است و نباید به‌عنوان رتبه‌بندی عمومی دو مدل یا تضمین رفتار آن‌ها در داده‌های دیگر تفسیر شود.
\end{enumerate}

\section{دستاوردهای اصلی}

\subsection{دستاوردهای علمی}

نخستین دستاورد علمی پژوهش، صورت‌بندی استدلال چندمرحله‌ای به‌عنوان منبعی مکمل به جای پاسخ قطعی سامانه است. کیفیت ظاهری ردپای استدلال صحت برچسب را تضمین نمی‌کند؛ ازاین‌رو، تولید تفسیر و انتخاب پاسخ دو وظیفۀ جدا هستند. دستاورد دوم، تفکیک بازبینی، کنترل میانی و انتخاب نهایی است. بازبین خطای احتمالی را ساختارمند می‌کند، کنترل‌گر تصمیمی قابل ردیابی می‌سازد و انتخاب‌گر با تکیه بر الگوهای آموزش منبع پاسخ را تعیین می‌کند.

دستاورد دیگر، ارزیابی انتخاب منبع بدون نشت اطلاعات آزمون است. پروفایل‌ها و آستانه‌های محافظ فقط با بخش آموزش ساخته و تنظیم و سپس پس از تثبیت روی آزمون اعمال شدند. بهبود حاصل نشان داد که می‌توان مکمل‌بودن خطاهای منابع را بدون دسترسی به پاسخ آزمون به سود عملی تبدیل کرد. در عین حال، فاصلۀ سامانۀ نهایی با اوراکل نشان می‌دهد که تشخیص منبع مناسب هنوز مسئله‌ای باز است.

\subsection{دستاوردهای فنی و مهندسی}

از نظر فنی، پروژه خط لوله‌ای انتهابه‌انتها از خواندن داده تا گزارش نهایی فراهم کرد. تقسیم رسمی آموزش و آزمون حفظ شد و خروجی‌های مستقیم، استدلالی، تشخیصی، کنترل‌گر و نهایی جداگانه ذخیره شدند. پیش از ارزیابی نیز شمار و یکتایی ردیف‌ها، هم‌ترازی شناسه، جمله، هدف و دامنه و اعتبار برچسب‌ها به‌صورت خودکار کنترل شد.

اجرای اصلی با \mbox{\lr{Qwen3 8B}} و \lr{Ollama} در یک چارچوب واحد انجام شد. برای مقایسۀ تکمیلی نیز خروجی دو مدل روی شناسه‌های مشترک سنجیده شد تا تفاوت جمعیت آزمون وارد مقایسه نشود. این اجزا امکان ردیابی خروجی هر مرحله و تکرار ارزیابی در محدودۀ تنظیمات گزارش‌شده را فراهم می‌کنند.

\section{محدودیت‌های پژوهش}

مجموعۀ اصلی شامل متن‌های انگلیسی و فقط دو دامنۀ لپ‌تاپ و رستوران است؛ بنابراین، تعمیم نتیجه به زبان فارسی، شبکه‌های اجتماعی یا حوزه‌های تخصصی به آزمایش مستقل نیاز دارد. نتیجۀ اصلی نیز به اجرای کامل \mbox{\lr{Qwen3 8B}} تعلق دارد. مقایسۀ \mbox{\lr{Gemini 2.5 Flash}} فقط روی ۹۰ نمونۀ آزمون مشترک انجام شد و شاهدی تکمیلی است.

مسیر استدلال چندمرحله‌ای غیرقطعی است. خودسازگاری بخشی از نوسان برچسب را مهار می‌کند، اما خطاهای مشترک اجراها را از بین نمی‌برد. ردپای استدلال نیز ارزیابی انسانی مستقلی از نظر صحت و وفاداری نداشته است. از سوی دیگر، مراحل چندگانه فراخوانی و متن بیشتری تولید می‌کنند و چون زمان‌سنجی کنترل‌شده انجام نشده، نسبت دقیق بهبود به زمان یا هزینه مشخص نیست.

اطلاعات تشخیصی در ساخت پروفایل‌ها نقش داشتند، اما برچسب پیشنهادی بازبین در هیچ نمونۀ آزمون منبع نهایی نبود؛ پس کارایی مستقل آن اثبات نشده است. برخی پروفایل‌ها کم‌نمونه‌اند، ولی سیاست اصلی برای پشتیبانی کمتر از ۱۰ یا حاشیۀ ناکافی به مسیر مستقیم بازمی‌گردد؛ بنابراین، خطر تصمیم‌های مبتنی بر گروه‌های کوچک محدود شده، نه اینکه کم‌نمونه‌بودن آن گروه‌ها حذف شده باشد. فاصلۀ سامانۀ نهایی با اوراکل فقط کرانی تحلیلی است و نباید عملکرد تضمین‌شدۀ انتخاب‌گر آینده تلقی شود.

در نهایت، مسیر \lr{THOR} برای مدل، داده و زیرساخت پروژه سازگار شده و بازتولید کامل تنظیم رسمی پژوهش اصلی نیست. وزن مدل‌ها ثابت مانده و اثر ریزتنظیم نیز بررسی نشده است. یافته‌های این پایان‌نامه در نتیجه به چارچوب پرامپت‌محور و انتخاب منبع محدود می‌شوند.

\section{کارهای آتی}

\subsection{توسعه و بهبود انتخاب‌گر منبع}

می‌توان شمار تقسیم‌های اعتبارسنجی را افزایش داد و پایداری انتخاب پروفایل‌ها و آستانه‌های محافظ را با بازنمونه‌گیری سنجید. هموارسازی الگوهای کم‌نمونه و ادغام پروفایل‌های مشابه نیز قابل بررسی‌اند. فرادسته‌بندهای کنترل‌شده می‌توانند با ویژگی‌های تنظیم‌شده از بخش آموزش و قید پرهیز از اصلاح کم‌اطمینان توسعه یابند. ارزیابی آن‌ها باید علاوه بر میانگین عملکرد، شمار پاسخ‌های درستی را که تضعیف می‌شوند گزارش کند و جدایی آموزش و آزمون را حفظ کند.

\subsection{بهبود بازبینی و استدلال هدف‌محور}

برای کاهش جابه‌جایی دامنۀ هدف، می‌توان بازۀ هدف و شواهد مرتبط را در تمام مراحل برجسته کرد و پیش از تعیین قطبیت، ارتباط شاهد با هدف را سنجید. آزمون سازگاری میان جنبه، نظر و برچسب نیز می‌تواند استدلال‌های ناهم‌جهت را آشکار کند. بازبین آینده بهتر است شواهد موافق و مخالف را تفکیک و نبود شاهد کافی را معتبر تلقی کند. اطمینان تشخیصی نیز باید با دادۀ اعتبارسنجی کالیبره و اثر تغییرها با کاهش خطاهای هدف‌محور سنجیده شود.

\subsection{تحلیل خطاهای مشترک و گسترش منابع}

تحلیل دقیق‌تر ۱۱۶ نمونۀ مشترک نادرست می‌تواند سهم دانش عمومی، ابهام برچسب، ساختار نحوی و چندهدف‌بودن جمله را روشن کند. برای این موارد باید منابعی با شیوۀ استنتاج یا دانش متفاوت، مانند پرامپت هدف‌محور مستقل یا بازیابی دانش زمینه‌ای، بررسی شوند. افزودن هر منبع باید با مطالعۀ حذف مؤلفه همراه باشد تا اثر تنوع واقعی از افزایش صرف تعداد فراخوانی‌ها جدا شود.

\subsection{گسترش ارزیابی به داده‌ها و مدل‌های دیگر}

ارزیابی روی مجموعه‌داده‌ها، دامنه‌ها و زبان‌های دیگر، به‌ویژه متن فارسی، می‌تواند پایداری سیاست انتخاب را روشن کند. مقایسۀ مدل‌ها نیز باید روی چند زیرمجموعۀ مستقل یا آزمون کامل انجام شود. در این بررسی، علاوه بر امتیاز نهایی، تغییر الگوهای خطا، مکمل‌بودن منابع و نیاز انتخاب‌گر به تنظیم دوباره با دادۀ آموزش اهمیت دارد.

\subsection{کاهش هزینۀ اجرا}

مراحل پرهزینه می‌توانند فقط هنگام نامطمئن‌بودن پیش‌بینی مستقیم یا وجود نشانۀ ابهام هدف فعال شوند. چنین دروازه‌ای باید کاهش فراخوانی را همراه با تغییر عملکرد گزارش کند. اندازه‌گیری کنترل‌شدۀ زمان، شمار توکن، حافظه و هزینۀ سرویس نیز امکان مقایسۀ پیکربندی‌ها برای کاربردهای برخط، دسته‌ای و کم‌منبع را فراهم می‌کند.

\subsection{ارزیابی انسانی و ریزتنظیم}

ردپای استدلال باید از نظر صحت شواهد، ارتباط با هدف، سازگاری با قطبیت و وفاداری ارزیابی انسانی شود. در صورت فراهم‌شدن دادۀ بزرگ‌تر و برچسب‌های معتبر برای جنبه، نظر، شواهد و نوع خطا، ریزتنظیم پارامترکم نیز قابل بررسی است. این مسیر باید با خط پایۀ پرامپت‌محور و تحت تقسیم یکسان مقایسه شود؛ زیرا آموزش با برچسب‌های خودتولیدشده ممکن است خطای مدل مولد را بازتولید کند.

\section{نتیجه‌گیری نهایی}

این پایان‌نامه نشان داد که افزودن مراحل بیشتر به یک پرامپت، به‌تنهایی تحلیل احساسات ضمنی را بهبود نمی‌دهد. پیش‌بینی مستقیم مرجعی قوی بود و استدلال چندمرحله‌ای، با وجود ایجاد ردپای قابل تحلیل، خطاهای تازه‌ای نیز وارد کرد. ارزش آن زمانی آشکار شد که به‌عنوان منبعی مکمل به کار رفت و انتخاب پاسخ از تولید استدلال جدا شد.

دستاورد اصلی پژوهش چارچوبی برای مدیریت اختلاف منابع است: خطای احتمالی ساختارمند می‌شود و پاسخ نهایی با سیاستی جداگانه و آموخته‌شده از آموزش انتخاب می‌گردد. بهبود مشاهده‌شده به داده و تنظیم حاضر مقید است، اما نشان می‌دهد توجه هم‌زمان به شواهد، محدودۀ هدف و اعتمادپذیری نسبی منابع مسیر مناسبی برای ادامۀ پژوهش است.
